% =====================================================================
\section{Fase de Pruebas del sistema}
\label{sec:pruebas}

La última fase de Mobile-D verifica que el producto cumpla sus
requerimientos. Las pruebas se organizaron en cinco niveles, de los más
automáticos a los que requieren al usuario (Tabla~\ref{tab:estrategia}).

\begin{table}[H]
\centering
\caption{Estrategia de pruebas}
\label{tab:estrategia}
\begin{tabularx}{\textwidth}{L{3.5cm} Y L{3.4cm}}
\toprule
\textbf{Nivel} & \textbf{Objeto de la prueba} & \textbf{Técnica} \\
\midrule
Verificación estática & Tipos, calidad y empaquetado de todo el código. &
Compilador y analizadores (sección~\ref{sec:estabilizacion}) \\
Lógica de negocio & Funciones puras: generador de rutinas y cálculo del
IMC. & Pruebas unitarias con valores límite \\
Reglas de seguridad & Aislamiento entre usuarios y validación de datos en el
servidor. & Escenarios permitidos y denegados \\
Funcionales & Cada requerimiento funcional en la aplicación instalada. &
Caja negra, casos de prueba \\
Aceptación & Cumplimiento de cada objetivo específico. & Trazabilidad
objetivo--requerimiento--prueba \\
\bottomrule
\end{tabularx}
\end{table}

\subsection{Pruebas de la lógica de negocio}

Las funciones \archivo{generarPlanHeuristico}, \archivo{repartirCantidad},
\archivo{calcularIMC} y \archivo{clasificarIMC} no dependen de la red ni de la
base de datos, por lo que se ejercitaron directamente con Node.js~22 sobre el
código fuente del proyecto, sin modificaciones en su lógica.

\paragraph{Generador de rutinas.} Se evaluaron las doce combinaciones posibles
de objetivo y nivel. La Tabla~\ref{tab:prueba-generador} muestra los
resultados obtenidos, que coinciden en todos los casos con las reglas de las
Tablas~\ref{tab:reglas-frecuencia} a~\ref{tab:reglas-division}.

\begin{table}[H]
\centering
\caption{Resultados del generador para las doce combinaciones de entrada}
\label{tab:prueba-generador}
\small
\begin{tabular}{L{2.5cm} L{2.2cm} C{0.9cm} C{1.9cm} L{4.4cm} C{1.1cm}}
\toprule
\textbf{Objetivo} & \textbf{Nivel} & \textbf{Días} & \textbf{Series × Rep.
(desc.)} & \textbf{División obtenida} & \textbf{Ejerc.} \\
\midrule
Ganar masa & Principiante & 3 & 4×9 (90 s) & Emp. -- Tir. -- Pie. & 15 \\
Ganar masa & Intermedio & 4 & 4×9 (90 s) & Emp. -- Tir. -- Pie. -- C.C. & 21 \\
Ganar masa & Avanzado & 5 & 4×9 (90 s) & Emp. -- Tir. -- Pie. -- Emp. -- Tir. & 25 \\
Mantenimiento & Principiante & 3 & 3×12 (60 s) & Emp. -- Tir. -- Pie. & 15 \\
Mantenimiento & Intermedio & 4 & 3×12 (60 s) & Emp. -- Tir. -- Pie. -- C.C. & 21 \\
Mantenimiento & Avanzado & 5 & 3×12 (60 s) & Emp. -- Tir. -- Pie. -- Emp. -- Tir. & 25 \\
Perder peso & Principiante & 3 & 3×15 (45 s) & Emp. -- Car. -- Pie. & 14 \\
Perder peso & Intermedio & 4 & 3×15 (45 s) & Emp. -- Tir. -- Car. -- Pie. & 19 \\
Perder peso & Avanzado & 5 & 3×15 (45 s) & Emp. -- Tir. -- Pie. -- Car. -- C.C. & 25 \\
Resistencia & Principiante & 3 & 3×18 (30 s) & Emp. -- Car. -- Pie. & 14 \\
Resistencia & Intermedio & 4 & 3×18 (30 s) & Emp. -- Tir. -- Car. -- Pie. & 19 \\
Resistencia & Avanzado & 5 & 3×18 (30 s) & Emp. -- Tir. -- Pie. -- Car. -- C.C. & 25 \\
\bottomrule
\end{tabular}
\fuente{Ejecución de \archivo{generarPlanHeuristico}. Emp.: empuje; Tir.:
tirón; Pie.: pierna; Car.: cardio; C.C.: cuerpo completo; Ejerc.: total de
ejercicios por semana.}
\end{table}

Además, \archivo{repartirCantidad} distribuyó correctamente los ejercicios
entre las categorías de cada bloque: cinco ejercicios en tres categorías se
reparten como 2--2--1, seis en tres como 2--2--2 y cuatro en dos como 2--2.

\paragraph{Cálculo del IMC.} Se probaron valores en los límites de cada
categoría de la OMS (Tabla~\ref{tab:prueba-imc}). Todos los casos produjeron
el resultado esperado.

\begin{table}[H]
\centering
\caption{Pruebas de valores límite del cálculo y la clasificación del IMC}
\label{tab:prueba-imc}
\begin{tabular}{C{2.2cm} C{2.4cm} C{2.2cm} L{3.3cm} C{2cm}}
\toprule
\textbf{Peso (kg)} & \textbf{Estatura (cm)} & \textbf{IMC} &
\textbf{Categoría obtenida} & \textbf{Resultado} \\
\midrule
50.0 & 170 & 17.3 & Bajo peso & Correcto \\
53.4 & 170 & 18.5 & Normal & Correcto \\
70.0 & 175 & 22.9 & Normal & Correcto \\
72.2 & 170 & 25.0 & Sobrepeso & Correcto \\
86.7 & 170 & 30.0 & Obesidad & Correcto \\
100.0 & 170 & 34.6 & Obesidad & Correcto \\
\midrule
\multicolumn{2}{c}{IMC = 18.4 / 18.5} & & Bajo peso / Normal & Correcto \\
\multicolumn{2}{c}{IMC = 24.9 / 25.0} & & Normal / Sobrepeso & Correcto \\
\multicolumn{2}{c}{IMC = 29.9 / 30.0} & & Sobrepeso / Obesidad & Correcto \\
\bottomrule
\end{tabular}
\fuente{Ejecución de \archivo{calcularIMC} y \archivo{clasificarIMC}.}
\end{table}

\subsection{Pruebas de las reglas de seguridad}

Las reglas de seguridad son la única barrera real frente a un acceso
malintencionado, porque las credenciales de Firebase están contenidas en la
aplicación. La Tabla~\ref{tab:prueba-reglas} define los escenarios evaluados
y el resultado que exigen las reglas del Anexo~\ref{anx:reglas}. Los
escenarios se ejecutan con el simulador de reglas de la consola de Firebase
(\emph{Rules Playground}), autenticando como un usuario $A$ y operando sobre
datos propios o de otro usuario $B$.

\begin{table}[H]
\centering
\caption{Escenarios de prueba de las reglas de seguridad}
\label{tab:prueba-reglas}
\small
\begin{tabularx}{\textwidth}{L{0.9cm} Y C{1.9cm} C{1.9cm}}
\toprule
\textbf{ID} & \textbf{Escenario} & \textbf{Esperado} & \textbf{Obtenido} \\
\midrule
SEG-01 & Usuario no autenticado lee \archivo{maquinas/\{id\}} & Denegado &
\pendiente{} \\
SEG-02 & $A$ lee \archivo{maquinas/\{id\}} & Permitido & \pendiente{} \\
SEG-03 & $A$ escribe en \archivo{maquinas/\{id\}} & Denegado & \pendiente{} \\
SEG-04 & $A$ lee \archivo{usuarios/A/rutinas} & Permitido & \pendiente{} \\
SEG-05 & $A$ lee \archivo{usuarios/B/rutinas} & Denegado & \pendiente{} \\
SEG-06 & $A$ crea una rutina con \archivo{diasPorSemana = 8} & Denegado &
\pendiente{} \\
SEG-07 & $A$ crea una rutina con \archivo{usuarioId = B} en su propia
subcolección & Denegado & \pendiente{} \\
SEG-08 & $A$ modifica el campo \archivo{email} de su perfil & Denegado &
\pendiente{} \\
SEG-09 & $A$ registra progreso con \archivo{peso = 0} & Denegado &
\pendiente{} \\
SEG-10 & $A$ registra ingesta con \archivo{calorias = -100} & Denegado &
\pendiente{} \\
SEG-11 & $A$ crea su perfil con \archivo{creadoEn} distinto de la hora del
servidor & Denegado & \pendiente{} \\
SEG-12 & $A$ sube un archivo a Storage & Denegado & \pendiente{} \\
\bottomrule
\end{tabularx}
\end{table}

\subsection{Pruebas funcionales}

Las pruebas funcionales se realizan con la técnica de caja negra sobre el
\emph{development build} instalado en un dispositivo Android, siguiendo los
criterios de aceptación definidos en cada iteración. Se definieron 30 casos de
prueba que cubren los doce requerimientos funcionales; la matriz completa, con
los pasos, el resultado esperado y el resultado obtenido, se encuentra en el
Anexo~\ref{anx:casos}. La Tabla~\ref{tab:resumen-funcionales} resume su
distribución.

\begin{table}[H]
\centering
\caption{Distribución de los casos de prueba funcionales}
\label{tab:resumen-funcionales}
\begin{tabular}{L{6cm} C{2.5cm} C{2.8cm}}
\toprule
\textbf{Módulo} & \textbf{Casos} & \textbf{Requerimientos} \\
\midrule
Autenticación, onboarding y perfil & 7 & RF1, RF2, RF12 \\
Guía de máquinas y ejercicios & 6 & RF3, RF4, RF5 \\
Escáner con realidad aumentada & 4 & RF6 \\
Rutinas y planificación & 6 & RF7, RF8 \\
Nutrición (IMC) & 3 & RF9 \\
Progreso y reportes & 4 & RF10, RF11 \\
\midrule
\textbf{Total} & \textbf{30} & \\
\bottomrule
\end{tabular}
\end{table}

\paragraph{Dispositivos de prueba.} \pendiente{indicar marca, modelo y versión
de Android de los dispositivos en que se ejecutaron las pruebas.}

\subsection{Pruebas de aceptación por objetivo}

Finalmente, la Tabla~\ref{tab:trazabilidad} relaciona cada objetivo
específico con los requerimientos que lo implementan y con las pruebas que lo
verifican. Esta trazabilidad es la base de la evaluación del cumplimiento de
objetivos que realiza el tribunal (Artículo~64 del Reglamento de Titulación).

\begin{table}[H]
\centering
\caption{Trazabilidad entre objetivos, requerimientos y pruebas}
\label{tab:trazabilidad}
\begin{tabularx}{\textwidth}{L{1.1cm} Y L{1.9cm} L{2.9cm}}
\toprule
\textbf{Obj.} & \textbf{Funcionalidad entregada} & \textbf{Req.} &
\textbf{Pruebas} \\
\midrule
OE1 & Guía de 35 máquinas con búsqueda, filtros, fichas y escáner con
realidad aumentada. & RF3, RF4, RF6 & CP-08 a CP-17 \\
OE2 & Ejercicios con animación, músculo objetivo y dificultad. & RF5 &
CP-12, CP-13 \\
OE3 & Onboarding de objetivo y nivel; generador de rutinas basado en
reglas. & RF2, RF7 & CP-04, CP-05, CP-18, CP-19, Tabla~\ref{tab:prueba-generador} \\
OE4 & Planificador de rutinas por día, rutina activa y agregado desde la
ficha. & RF8 & CP-20 a CP-23 \\
OE5 & Calculadora de IMC con clasificación e historial. & RF9 & CP-24 a
CP-26, Tabla~\ref{tab:prueba-imc} \\
OE6 & Registro de peso, estatura e IMC. & RF10 & CP-27, CP-28 \\
OE7 & Reporte con gráfico, variación y categoría actual. & RF11 & CP-29,
CP-30 \\
\bottomrule
\end{tabularx}
\end{table}
